\documentclass[10pt,a4paper]{amsart}
\usepackage[margin=19mm]{geometry}
\usepackage{amsmath,amssymb,mathtools}
\usepackage[scr=rsfs,cal=euler]{mathalfa}
\usepackage{xfrac}
\usepackage[unicode,hidelinks,bookmarks=false]{hyperref}
\newcommand{\ie}{i.e.\ }
\newcommand{\cf}{cf.\ }
\newcommand{\resp}{respectively\ }
\newcommand{\Subsection}{\subsection}
\providecommand{\nobreakdash}{\nobreak\hbox{-}}
\setlength{\emergencystretch}{2em}
\multlinegap0pt
\usepackage{bm}
\usepackage{bbm}
\usepackage{dsfont}
\usepackage{xspace}
\usepackage{cancel}
\usepackage{mathtools}
\usepackage[shortlabels]{enumitem}
\usepackage{amsthm}
\makeatletter
\@ifundefined{theorem}{\newtheorem{theorem}{Theorem}}{}
\makeatother
\newcommand{\mc}{\mathcal}
\newcommand{\on}{\operatorname}
\title[Lie symmetry analysis of the two-Higgs-doublet model field e]{Lie symmetry analysis of the two-Higgs-doublet model field equations}
\author{M. Aa. Solberg}
\date{Source: arXiv:2510.09542v2 (CC BY 4.0)}
\begin{document}
\maketitle

\noindent\emph{Source and licence.} Lie symmetry analysis of the two-Higgs-doublet model field equations. Version arXiv:2510.09542v2 (CC BY 4.0), distributed under the Creative Commons Attribution 4.0 International licence, as recorded in the arXiv licence metadata for this exact version and shown on the abstract page https://arxiv.org/abs/2510.09542v2. Full rights notice: licenses/arxiv-2510.09542-CC-BY-4.0.txt.\par\smallskip

\noindent\textit{Editorial scope.} Counted unit: the Theorem whose source label is \texttt{P:prXannTannV} in the article (source0.tex lines 568--585, source label \texttt{P:prXannTannV}), delivered together with the article's own complete original proof of it (source0.tex lines 586--626, one \texttt{proof} environment of 41 lines). No mathematics is written, repaired or re-derived: statement and proof are the article's own text line for line. Required context retained uncounted: E:prDelta=0PDEsymCond (lines 370--372); E:condStrictVarSym (lines 439--441); E:EtotDiv (lines 471--473); E:divSymmetry (lines 475--477); E:adjFrechetScAff (lines 533--535); E:commFrechet (lines 540--542). Every definition, display and label of those lines is retained unchanged. Omitted from this package and not redistributed: 1--369, 373--438, 442--470, 474--474, 478--532, 536--539, 543--567, 627--1738. Nothing inside a retained excerpt is omitted or altered: every retained body is delivered exactly as the article's lines are written, so no omission receipt of any kind accompanies this package. Citations: the retained excerpts carry no citation command at all, so no bibliography entry of the article is transcribed and no reference is redistributed. Reference adaptation: none was needed and none was made; every reference inside the delivered unit resolves inside it, so no unresolved reference or citation is produced. Printed slips kept literal: no printed word, symbol, hypothesis or number of the counted statement or of its proof was altered. Layout and class: the article's own class is \texttt{article} with packages including inputenc, a4wide, amsmath, amsfonts, amssymb, amsthm, mathabx, array, subfigure, cite, xcolor, soul, cancel, appendix; the wrapper is the lane's pinned \texttt{amsart} editorial wrapper at 10pt on A4, which restates the theorem-like environments the retained text needs and adds the article's own macro definitions that the retained text needs (\texttt{\textbackslash mc}, \texttt{\textbackslash on}), with extra wrapper packages (bm, bbm, dsfont, xspace, cancel, mathtools, enumitem). The delivered numbering is the unit's own. Assets: no retained span contains a figure environment, a graphics-inclusion command, a PDF-page inclusion, a table or a TikZ picture; no image file is copied into this package, the package declares no asset dependency and no image file is redistributed with it. Licence: version arXiv:2510.09542v2 of this article is distributed under the Creative Commons Attribution 4.0 International licence, as recorded in the arXiv licence metadata for this exact version and shown on the abstract page https://arxiv.org/abs/2510.09542v2. A full rights notice is delivered at licenses/arxiv-2510.09542-CC-BY-4.0.txt. Characters: every retained excerpt is pure ASCII, so the delivered engine, XeLaTeX with its default Latin Modern text font, reports no missing character.
\begin{align}\label{E:prDelta=0PDEsymCond}
	 (\on{pr}X(\Delta_i))|_{\Delta = 0}=0 \quad \forall i\in \{1,\ldots,m\},  
\end{align}

\begin{align}\label{E:condStrictVarSym}
	\on{pr}X(\mc{L})+\mc{L}\,{d_\mu \xi^\mu}=0,
\end{align}

\begin{align}\label{E:EtotDiv}
	E(f)=0.
\end{align}

\begin{align}\label{E:divSymmetry}
	\on{pr}X(\mc{L})+\mc{L}\,d_\mu \xi^\mu=f,
\end{align}

\begin{align}\label{E:adjFrechetScAff}
	(D_Q)_{ij}^\ast=  \frac{\partial \eta^j(y^1,\ldots,y^{q})}{\partial y^i}=\mc{J}_{ji}=\mc{J}^T_{ij},
\end{align}

\begin{align}\label{E:commFrechet}
	E(\on{pr}X_Q(\mc{L}))=\on{pr}X_Q(E(\mc{L}))+D_Q^\ast E(\mc{L}),
\end{align}

\begin{theorem}\label{P:prXannTannV}
Let $\mc{L}=T-V$ be a polynomial potential theory and let the infinitesimal generator 
\[
X=\eta^i(y^1,\ldots,y^{q})\partial_{y^i}, 
\]
where each $\eta^i$ is a polynomial, be a symmetry of $E(\mc{L})=0$. Moreover, assume that either $V(\varphi_1,\ldots,\varphi_m)$ does not contain any linear terms $\alpha_i \varphi_i$, or that
 $\alpha_i\ne 0 \Rightarrow \eta^i(y^1,\ldots,y^{q})$ does not
contain a constant term.
%consider the action of the prolongation of $X$
%on the kinetic terms $T$ of $\mc{L}$,
%$\on{pr}X(T)$. 
Then,   
\begin{enumerate}[label=(\roman*)]
 \item If $\on{pr}X(T)=0$, the symmetry generated by $X$ is strictly variational.
  \item If $\on{pr}X(T)=d_\mu \beta^\mu$ for some non-vanishing current $\beta^\mu$, i.e.~a total divergence, 
        the symmetry generated by $X$ is a divergence symmetry.
\end{enumerate}
\end{theorem}

\begin{proof}
 The linearized symmetry condition \eqref{E:prDelta=0PDEsymCond}
applied on the Euler-Lagrange equations $E(\mc{L})=0$ yields
\begin{align}\label{E:prDelta=0PDEsymCond2}
	 \big(\on{pr}X(E_i(\mc{L}))\big)|_{E(\mc{L}) = 0}=0, \quad \forall i\in \{1,\ldots,q\},  
\end{align}
while \eqref{E:commFrechet} implies
\begin{align}\label{E:commFrechet2}
	\on{pr}X(E(\mc{L}))= E(\on{pr}X(\mc{L}))-D_Q^\ast E(\mc{L}),
\end{align}
since $\on{pr}X= \on{pr}X_Q$, because the former is already in evolutionary form. Assume 
\begin{align}
\on{pr}X(T)=d_\mu \beta^\mu, 	
\end{align}
for a possible vanishing current $\beta$. Substituting \eqref{E:adjFrechetScAff} into \eqref{E:commFrechet2} then yields, 
\begin{align}\label{E:commFrechet3}
	\on{pr}X(E(\mc{L}))= -E(\on{pr}X(V))-\mc{J}^T E(\mc{L}),
\end{align}
because $E(d_\mu \beta^\mu)=0$ cf.~\eqref{E:EtotDiv}, and
where $\mc{J}$ is the Jacobian matrix of $\eta$.
Then, if $E(\mc{L}) = 0$, the first and last terms in \eqref{E:commFrechet3} vanish, cf.~\eqref{E:prDelta=0PDEsymCond2}, that is, for any $i$
\begin{align}\label{E:derPrPot0}
	E_i(\on{pr}X(V))=\frac{\partial}{\partial y^i}\on{pr}X(V)=0,
\end{align}
because there are no derivatives in the potential. The polynomial equation \eqref{E:derPrPot0} will also hold when $E(\mc{L})\ne 0$, because $E(\mc{L})= 0$ does not imply any polynomial relations (i.e.,~consequences) between the fields of $V$, as $\mc{L}$ is a polynomial potential theory. But then $\on{pr}X(V)=C$ for constant $C$, and $C=0$ because all terms of 
\begin{align}
\on{pr}X(V)=X(V)=\eta^i\partial_{y^i}V	
\end{align}
are, if non-vanishing, at least linear in the fields, because if $\partial_{y^i}V$ includes a constant term for an $i$, then $\eta^i$ does not, per assumption. This means that the potential $V$ is annihilated by 
the prolongation of $X$,
\begin{align}
	\on{pr}X(V)=0, 
\end{align}
  and hence 
	\begin{align}
		\on{pr}X(\mc{L})=\on{pr}X(T)=d_\mu \beta^\mu.
	\end{align}
Thus, if 
 $\beta=0$ then $X$ generates a strict variational symmetry, cf.~\eqref{E:condStrictVarSym} with $\xi=0$.
Furthermore, if $\beta$ is non-vanishing, $X$ generates a divergence symmetry, cf.~\eqref{E:divSymmetry} with $\xi=0$.
\end{proof}

\end{document}
